\chapter{Tablas de conversión}
\label{ap:conversiones}

\fuente{POH, sección 1}

\begin{table}[htbp]
\centering\small
\caption{Sistema métrico a sistema inglés.}
\begin{tabular}{@{}lrl@{}}
\toprule
Multiplique & por & para obtener \\
\midrule
centímetros (cm) & 0,3937 & pulgadas (in) \\
kilogramos (kg) & 2,2046 & libras (lb) \\
kilómetros (km) & 0,5400 & millas náuticas (MN) \\
kilómetros (km) & 0,6214 & millas terrestres (mi) \\
litros (l) & 0,2642 & galones EE.\,UU. (gal) \\
litros (l) & 1,0567 & cuartos (qt) \\
metros (m) & 3,2808 & pies (ft) \\
milibares (mb) & 0,0295 & pulgadas de mercurio (inHg) \\
\bottomrule
\end{tabular}
\end{table}

\begin{table}[htbp]
\centering\small
\caption{Sistema inglés a sistema métrico.}
\begin{tabular}{@{}lrl@{}}
\toprule
Multiplique & por & para obtener \\
\midrule
pies (ft) & 0,3048 & metros (m) \\
galones EE.\,UU. (gal) & 3,7854 & litros (l) \\
pulgadas (in) & 2,5400 & centímetros (cm) \\
pulgadas (in) & 25,4000 & milímetros (mm) \\
pulgadas de mercurio (inHg) & 33,8638 & milibares (mb) \\
millas náuticas (MN) & 1,8520 & kilómetros (km) \\
libras (lb) & 0,4536 & kilogramos (kg) \\
cuartos (qt) & 0,9464 & litros (l) \\
millas terrestres (mi) & 1,6093 & kilómetros (km) \\
\bottomrule
\end{tabular}
\end{table}

\begin{itemize}
\item 1 milla náutica = 1,1508 millas terrestres; 1 milla terrestre = 0,8690 millas náuticas.
\item $^\circ\mathrm{F} = \tfrac{9}{5}\,(^\circ\mathrm{C}) + 32$; \quad $^\circ\mathrm{C} = \tfrac{5}{9}\,(^\circ\mathrm{F} - 32)$.
\end{itemize}

\section*{Conversiones útiles adicionales}

Estas equivalencias no figuran en el POH; se agregan para el trabajo de taller.

\begin{itemize}
\item Torque: 1~in-lb = 0,113~N·m; 1~ft-lb = 12~in-lb = 1,356~N·m.
\item Presión: 1~psi = 6,895~kPa = 0,0689~bar.
\item Combustible de aviación (AVGAS): 6~lb/gal, valor usado por el POH para el peso y balance (aprox. 0,72~kg/l).
\item RPM: 100\,\% en el tacómetro = 2550~RPM de motor = 510~RPM de rotor; 104\,\% = 2652 / 530~RPM.
\end{itemize}
